\section{Problem and Method}\label{sec:method}
\subsection{Setup and notation}
The input is a matrix $X \in \mathbb{R}^{n \times g}$ after preprocessing, with $n$
cells and $g$ genes. A clustering method takes $X$ and one setting $s$ and returns a
partition $P_s$ of the cells, with $|P_s|$ clusters. For MICA in small-data mode the
setting is the number of clusters $k \in \{2,\dots,20\}$, so $|P_s| = k$. For Leiden and
for MICA in graph mode it is a resolution $r > 0$, and $|P_s|$ is whatever community
detection returns, rising with $r$ for Leiden and falling with it in MICA's Louvain. Reference labels $y \in \{1,\dots,c\}^n$ exist for
every dataset, where $c$ is the number of annotated cell types, and are used only for
scoring.

\subsection{MICA, the method under test}
MICA is the clustering module of scMINER \cite{pan2025scminer}. We run the authors'
released implementation at commit \texttt{54fae81}, not a reimplementation. It scores a
pair of cells by binned mutual information. Each cell is discretised into
$b = \lfloor g^{1/3} \rfloor$ equal-width bins between that cell's own minimum and maximum
value, giving a discrete variable $C_i$ per cell. MICA turns the mutual information between two
cells into a distance in two ways. Graph mode uses
\begin{equation}
D(i,j) \;=\; 1 - \frac{I(C_i; C_j)}{H(C_i, C_j)},
\label{eq:mi}
\end{equation}
and small-data mode uses
\begin{equation}
D_{\mathrm{sd}}(i,j) \;=\; 1 - \frac{I(C_i; C_j)}{\sqrt{H(C_i)\,H(C_j)}},
\label{eq:misd}
\end{equation}
where $I$ is mutual information and $H$ entropy, in nats, estimated from the joint histogram of
the two cells' bin codes across the $g$ genes. Both lie in $[0,1]$. Dividing by the joint entropy
makes Eq.~\ref{eq:mi} a metric \cite{vinh2010information}, which Eq.~\ref{eq:misd} is not
guaranteed to be, and the two modes also round the bin edges slightly differently.

MICA then takes one of two paths, switching at 5,000 cells. Below that, \emph{small-data
mode} applies classical multidimensional scaling to the matrix of Eq.~\ref{eq:misd}, keeping 19
dimensions, and runs $k$-means. Above it, \emph{graph mode} builds an 80 nearest neighbour graph
under Eq.~\ref{eq:mi} with HNSW
\cite{malkov2020efficient}, embeds that graph with node2vec random walks, and runs Louvain
over a resolution grid. Nine of our ten datasets fall in small-data mode, which is why the
setting is $k$ for most results. Our own implementation of both forms agrees with the authors'
code to within $10^{-6}$ on Zeisel, where a test enforces it, and within $10^{-4}$ on the other
inputs we compared. It is used only where the released code offers no entry point.

\subsection{The baseline}
The baseline is the pipeline most practitioners run: principal component analysis to 50
components, a 15 nearest neighbour graph on those components, and Leiden community
detection \cite{traag2019louvain} through Scanpy. Its setting is the resolution.

\subsection{The selection problem}
The published benchmark fixes the setting with the labels: every method is either given
the true number of types $c$, or has its resolution swept and keeps only runs with
$|P_s| = c$, breaking ties by agreement with $y$. That is a well-posed question, but not
the one a user faces. We therefore require a \emph{selection rule}: a function mapping the
candidate set $\{P_s\}$ to one setting $\hat{s}$ using $X$ alone. The same rule is applied
to both methods, and only afterwards is $P_{\hat{s}}$ scored against $y$.

\subsection{The selection rules}
\textbf{Default.} The setting the software ships with: resolution 1.0 for Scanpy's
Leiden, 1.822 for MICA's graph mode, scored at 1.8, the nearest point on its grid, and
four clusters for MICA's small-data mode. These are hard-coded with no stated
derivation, so if they were chosen against these same benchmark datasets they carry
label information. We treat this as a weak rule for that reason.

\textbf{Silhouette.} Keep the setting with the highest mean silhouette. In Assignment 2
this was computed in each method's own embedding, which confounds the partition with the
space. Here it is computed for both methods in one shared space, the 50 principal
components, so that only the partitions differ.

\textbf{Stability.} Rerun the method on $B$ random subsamples of $80\%$ of the cells and
prefer settings whose clusters survive resampling, following chooseR
\cite{pattersoncross2021chooser}. For setting $s$, let $f_{ij}$ be the fraction of
subsamples containing both cells $i$ and $j$ that place them in the same cluster. Treating
$1 - f_{ij}$ as a distance, the score $\sigma_s$ is the median across clusters of the mean
within-cluster silhouette taken against the full-data partition $P_s$. A bootstrap over
cells gives an interval for $\sigma_s$. The threshold is the largest lower bound over all
settings, and chooseR keeps the setting with the most clusters clearing it
(Appendix~\ref{app:repro} gives the estimator).

That last step encodes a preference for fine partitions. Because reproducibility says
nothing about whether a split is real, we also report the plain maximum,
$\hat{s} = \arg\max_s \sigma_s$ with ties going to fewer clusters. The two readings share
all the computation. We call them \emph{finest stable} and \emph{most stable}.

\subsection{An exactness property that makes the rule affordable}
For MICA, rerunning on a subsample would appear to need a new distance matrix, the
dominant cost. It does not. In Eq.~\ref{eq:mi} cell $i$'s bins come from its own range and
the histogram runs over genes, so no term depends on a third cell, and the matrix
restricted to a subset equals the matrix computed from that subset alone. We compute it
once per dataset and take submatrices, which a test checks rather than assumes
(Appendix~\ref{app:exact}). Every later step, the multidimensional scaling included, is
redone per subsample.
